\part{Deletion frames}\label{part:frames}
\section{Marked-index frames and transverse tensor projections}
\label{ten:frames}
This section proves a uniform unweighted Schatten bound. Two complementary row-column estimates are needed. One resolves functions by subsets of marked indices and uses cancellation of isolated contacts; the other realizes a diagram in a signed tensor space and normalizes row densities before comparing them with columns.
\smallskip
All notation introduced below is local to this section. Traces are unnormalized, logarithms are natural, and a sweep on $2^d$ positions takes $d$ physical shuffles. The representation-theoretic conventions are those of Section~\ref{sec:conventions}.

\subsection{Two transverse product estimates}
\label{ten:frames:part:1}
Write $n=2^d$ in this section, and write $\mu_d=K_d$ for the sweep. The two product estimates below measure different features of the same representation: its first-row deficit, and its size outside a small hook. Neither alone supplies the error needed at every shape. Their combination gives the following unweighted conclusion. Here $f^\lambda=D_\lambda$, and subscripts denote Fourier matrices.

\begin{proposition}[Uniform Schatten estimate]
\label{ten:frames:main}
There is an absolute real exponent $8\le p<\infty$ such that, for every $d\ge1$ and every $\lambda\vdash n$ with $n=2^d$,
\begin{equation}
 f^\lambda \|(\mu_d)_\lambda\|_p^p \ \le\ 1\qquad (\lambda\vdash n).
 \label{ten:frames:eq:1}
\end{equation}

\end{proposition}

We use the branching rule \cite[Theorem~2.8.3]{Sagan2001}, Pieri's rule \cite[Theorem~7.15.7]{Stanley1999}, and the Littlewood--Richardson rule \cite[Theorem~16.4]{James1978}, together with the hook-length formula \cite[Theorem~3.10.2]{Sagan2001}. Ordinary Schur--Weyl duality \cite[Theorem~4.57 and Corollary~4.59]{etingof} supplies the tensor realizations. Recall the following consequences and conventions:
\begin{itemize}
\item The space of functions on injections of \(j\) distinct indices into \(h\) positions carries the representation induced from the trivial representation of \(\mathfrak S_{h-j}\). It contains only shapes with first row at least \(h-j\). For \(\lambda_1=h-j\), the multiplicity is \(f^{\lambda_*}\), where \(\lambda_*\) is the shape after the first row. In this case
\begin{equation}
 e^{-j}\binom hj f^{\lambda_*}\ \le\ f^\lambda\ \le\ \binom hj f^{\lambda_*}.
 \label{ten:frames:eq:2}
\end{equation}
Indeed the hook correction to the first-row factorial has log at most \(j\).
\item For polynomial irreducibles of the unitary group in \(s\) dimensions, indexed by partitions \(u\) of degree \(\le n\) with at most \(s\) rows, we use dimension \(\le (n+s+1)^{O(s^2)}\) and the highest weight \(u\). Thus a positive semidefinite matrix with ordered eigenvalues \(x_i\) decreasing acts (by polynomial extension) with operator norm \(\prod_i x_i^{u_i}\), which is also a lower bound for the trace of that action. The Schur character and dimension formulas are given in \citet[Theorem~4.63]{etingof}; the largest monomial follows from the tableau description and the highest weight.
\item All group algebra blocks can be prescribed separately. Projections selecting a subspace in one irreducible block and zero elsewhere are orthogonal projections in every unitary representation. For a transitive permutation representation (uniform probability normalization on the orbit), the squared norm of evaluation at a point on the subspace selected by such a projection \(B\) in type \(\alpha\) is
\[
 f^\alpha \operatorname{Tr}(B Q_\alpha),
\]
where \(Q\) is the average on the stabilizer of the point. This follows by matrix-entry orthogonality or the regular representation formula, realizing functions on cosets inside the regular representation.
\end{itemize}

Constants and big-O in estimates below are absolute, and logs in estimates can be taken natural. We give some matrix product estimates for the induction. Split the bits of a large sweep into two consecutive groups of nearly equal sizes. This gives a board of \(b\) rows and \(m\) columns, \(bm=n\), both within a constant factor of \(\sqrt n\). The two parts of the sweep operate along the rows and along the columns, respectively; write
\[
 H=\mathfrak S_m^b,\qquad K=\mathfrak S_b^m
\]
for the corresponding groups (the estimates work in either product order). The logarithms of the numbers of types (irreducibles) for \(H,K\) are \(O(n^{3/4}\log n)\), by standard partition bounds, or simply by splitting each partition diagram according to a square of side about the square root of its size. We denote the dimension of a subgroup type by \(h_A\) when the type is \(A\).

Here are the two estimates. Suppose \(X\in\mathbb C H,\ Y\in\mathbb C K\) satisfy
\begin{equation}
 h_A\|X_A\|_v^v\le 1,\qquad h_B\|Y_B\|_v^v\le 1
 \label{ten:frames:eq:3}
\end{equation}
on all their types. Put \(j=n-\lambda_1\), \(j>0\).

\begin{itemize}
\item With \(v=8\), we have
\begin{equation}
 f^\lambda\|(XY)_\lambda\|_8^8 \ \le\ \exp\{O(j+n^{0.8})\}.
 \label{ten:frames:eq:4}
\end{equation}
\item Put \(s=\lfloor n^{0.01}\rfloor\), and let \(k\) be the number of boxes of \(\lambda\) outside its first \(s\) rows and first \(s\) columns (outside their union). With \(v=2\), we have
\begin{equation}
 f^\lambda\|(XY)_\lambda\|_2^2
 \ \le\ \binom nk \exp\{O(k+n^{0.8})\}.
 \label{ten:frames:eq:5}
\end{equation}
Bounds with adjoints or reversed order follow the same way. We prove the estimates in detail.
\end{itemize}

\subsection{A deletion frame for injection spaces}
\label{ten:frames:part:57}
The selected-rank estimate will resolve a row-type vector into functions that remember only a prescribed number of marked positions in each row. The following lemma gives a bounded resolution before any row-column comparison is made.

\begin{lemma}[Deletion frame]\label{lem:deletion-frame}
Let $m\ge1$ and $0\le t\le q\le m$ be integers, and let $\alpha\vdash m$ have first row $m-t$. Equip the space $\mathcal J_q$ of functions on injections $[q]\to[m]$ with uniform probability measure. For $U\subseteq[q]$, let $E_U^{(q)}$ be conditional expectation retaining the coordinates in $U$. On the $\alpha$-isotypic subspace,
\[
 F_{q,t}:=\sum_{\substack{U\subseteq[q]\\|U|=t}}E_U^{(q)}
 \succeq e^{-Cq}I
\]
for an absolute constant $C$.
\end{lemma}
\begin{proof}
We compare successive numbers of indices. For $t<r\le q$, put
$B_r=\sum_{i=1}^r E_{[r]\setminus\{i\}}^{(r)}$, and let $\beta_r$ be its least eigenvalue on type $\alpha$. Pullback from injections indexed by $[r]\setminus\{i\}$ is an isometry $J_i$ into $\mathcal J_r$. Its adjoint is conditional expectation, so $J_iJ_i^*=E_{[r]\setminus\{i\}}^{(r)}$. The tower property gives
\[
 \sum_{i=1}^r J_iF_{r-1,t}^{(i)}J_i^*=(r-t)F_{r,t}:
\]
each $t$-subset is counted for exactly $r-t$ omitted indices. All these maps are $S_m$-equivariant. Starting with $F_{t,t}=I$, a lower bound $a_{r-1}$ at level $r-1$ therefore gives $a_r=a_{r-1}\beta_r/(r-t)$ at level $r$. It remains to bound
\begin{equation}\label{eq:deletion-product}
 \prod_{r=t+1}^q\frac{\beta_r}{r-t}\ge e^{-Cq}.
\end{equation}

Write $c_\zeta$ for the sum of the box contents of a partition $\zeta$. The multiplicity space of $\alpha$ in $\mathcal J_r$ is its $S_{m-r}$-invariant space, up to duality. Its $S_r$-types are the $\delta\vdash r$ for which $\alpha/\delta$ is a horizontal strip, by Pieri. On such a type the eigenvalue of $B_r$ is
\[
 \frac{r+c_\alpha-c_\delta-c_{(m-r)}}{m-r+1}.
\]
Indeed, a deletion averages the subgroup allowing one retained point to permute with the $m-r$ unretained points. On the invariant space this is identity plus the transpositions involving that point, divided by $m-r+1$. Summing the mixed transpositions subtracts the transposition sums for $S_r$ and $S_{m-r}$ from that for $S_m$. The sum on a type $\zeta$ acts by $c_\zeta$: sum the Jucys--Murphy elements and their content eigenvalues in \citet[Section~2, Eq.~(2.1), and Theorem~5.8]{VershikOkounkov2005}. This gives the formula.

To estimate these eigenvalues put $M=m-r$ and $a=m-t$. The $a$ nonempty columns of $\alpha$ form plateaus of equal height. A horizontal-strip removal uses a bottom box in each of $M$ distinct columns and must use a suffix within each plateau. For removed column indices $j_1<\cdots<j_M$, set $S_M=\sum_{l=1}^M(j_l-l)$. The eigenvalue numerator is
\[
 r+S_M-\sum_{\text{removed columns}}(\text{height}-1)
 \ge r-t+S_M.
\]
The smallest sum of removed column indices fills the earliest plateaus first and then uses a suffix in one partial plateau: moving a removal from a later plateau to the next legal column of an earlier unfilled plateau decreases its index. If $M=L+x$, where that partial plateau has width $w$ and follows $L$ columns, then $S_M\ge x(w-x)$. Hence
\begin{equation}\label{eq:deletion-ratio}
 \frac{\beta_r}{r-t}
 \ge\frac{a-M+x(w-x)}{(M+1)(a-M)}
 \ge\frac{(x+1)(w-x)}{(a+1)^2}.
\end{equation}
The last inequality uses $a-M\ge w-x$.

First suppose $q\ge m/4$. The last lower bounds in \eqref{eq:deletion-ratio} are at most one, so their product over the required interval $m-q\le M<m-t$ is at least their product over all $0\le M<a$. This latter product is
\[
 \prod_{\text{plateaus}}\frac{(w!)^2}{(a+1)^{2w}},
\]
whose negative logarithm is at most $O(a)+2\sum_w w\log(a/w)$. The distribution assigning mass $w/a$ to each plateau height has mean $m/a$. Comparison with a geometric distribution bounds its entropy by $O(1)+\log(m/a)$. Thus the logarithmic loss is $O(a+a\log(m/a))=O(m)=O(q)$.

If $q<m/4$, at most $t\le q$ columns have height greater than one. Every $M$ in the required interval lies in the final height-one plateau. Here $w-x=a-M$ and $x\ge m-q-t$. The first bound in \eqref{eq:deletion-ratio} gives $(x+1)/(M+1)\ge1/2$, so at most $q$ factors again cost $e^{O(q)}$. This proves \eqref{eq:deletion-product}. If $t=q$ the product is empty; for $t=0$ the type is trivial and all its ratios equal one.
\end{proof}

\subsection{Projection overlap on marked indices}
\label{ten:frames:part:49}
We now use the deletion frame on each row and column. The selected ranks below count directions inside subgroup carriers; each ambient projection includes every multiplicity copy.
\begin{lemma}[Selected-rank projection overlap]\label{lem:selected-rank}
Let $n=bm$, $\lambda\vdash n$, and $j=n-\lambda_1>0$. Let $H=S_m^b$ and $K=S_b^m$ be the row and column groups. Let $P_R,P_C$ be group-algebra orthogonal projections supported on one type each, whose carrier dimensions are $h_R,h_C$ and whose selected ranks are $r_R,r_C$. Then
\begin{equation}
 f^\lambda\|(P_R P_C)_\lambda\|_8^8
 \le e^{O(j)}h_Rr_Rh_Cr_C.
 \label{ten:frames:eq:6}
\end{equation}
\end{lemma}
\begin{proof}
If either selected rank is zero, its group-algebra action is zero and the result is immediate. Assume henceforth that $1\le r_R\le h_R$ and $1\le r_C\le h_C$.

Let $\mathcal B$ be the space of functions on $([b]\times[m])^j$ with uniform product measure, and let $\mathcal I\subseteq\mathcal B$ consist of functions supported on injections. Write $D(x)$ for the distinctness indicator. The $\lambda$-isotypic part $\mathcal I_\lambda$ consists of $f^{\lambda_*}$ copies of $V_\lambda$. Averaging any one board coordinate kills this part: functions using only the other $j-1$ coordinates have first row at least $n-j+1$, as follows by splitting them according to their equality patterns. Thus the product-space projection
\[
 Z=\prod_{i=1}^j(I-E_i),
\]
where $E_i$ averages coordinate $i$, fixes $\mathcal I_\lambda$.

\paragraph{Marked-function synthesis.}
Write the row type as $A=(\alpha_i)_{i=1}^b$, put $t_i=m-\alpha_{i,1}$, and let $t_R=\sum_i t_i$. If either selected type is absent from $\mathcal I$, the overlap vanishes. We may therefore assume both occur, so $t_R\le j$. For $U\subseteq[j]$ with $|U|=t_R$, let $\mathcal F_{R,U}$ consist of functions
\begin{equation}
 F\big((\operatorname{row}(x_i))_{i\le j},
       (\operatorname{col}(x_i))_{i\in U}\big)
 \label{ten:frames:eq:7}
\end{equation}
with product-measure norm. On a row assignment having $q_i$ indices in row $i$, such a function vanishes unless $q_i\le m$, $U$ marks exactly $t_i$ indices in that row, and the marked columns are distinct. Its marked-column dependence is selected by $P_R$.

Define $A_{R,U}F=DF$, and regard the synthesis map
$A_R(F_U)=\sum_U A_{R,U}F_U$ as a map into the ambient space $\mathcal B$, with range in $\mathcal V_R=\operatorname{ran}(P_R|_{\mathcal I})$. All adjoints below use this ambient codomain. On a feasible row-assignment fiber, let $E_U^{\rm inj}$ retain the marked columns by conditional expectation under the uniform injection law, and set it to zero when $U$ has the wrong counts. On $\mathcal V_R$,
\[
 A_{R,U}A_{R,U}^*=c(q)E_U^{\rm inj},\qquad
 c(q)=\prod_i\frac{(m-t_i)_{q_i-t_i}}{m^{q_i-t_i}}\ge e^{-O(j)}.
\]
Here $(m)_t$ is the falling factorial. The factor $c(q)$ is the probability that the unmarked columns complete the marked ones to injections; the bound follows from $(m)_t/m^t\ge e^{-O(t)}$. Conditional expectation commutes with the row group and hence with arbitrary carrier selections in $P_R$. Applying Lemma~\ref{lem:deletion-frame} in each row therefore gives $A_RA_R^*\succeq e^{-O(j)}I$ on $\mathcal V_R$. Its restriction there is invertible, and
\[
 S_R=A_R^*\big((A_RA_R^*)|_{\mathcal V_R}\big)^{-1},\qquad
 A_RS_R=I_{\mathcal V_R},\qquad \|S_R\|\le e^{O(j)}.
\]
Construct $\mathcal F_{C,V},A_C,S_C$ analogously for column marks $V$ of size $t_C$.

Let $I_R,I_C$ include $\mathcal V_R\cap\mathcal I_\lambda$ and $\mathcal V_C\cap\mathcal I_\lambda$ into $\mathcal B$. Restrict the domains of $S_R,S_C$ to these spaces. Then $A_RS_R=I_R$ and $A_CS_C=I_C$, and the overlap has the exact factorization
\begin{equation}\label{eq:frame-factorization}
 I_R^*I_C=S_R^*\mathbf L S_C,\qquad
 \mathbf L=A_R^*ZA_C=(L_{U,V})_{U,V},\qquad
 L_{U,V}=A_{R,U}^*ZA_{C,V}.
\end{equation}
The insertion of $Z$ uses the reconstructed vectors in $\mathcal I_\lambda$; individual summands $DF$ need not be fixed by $Z$. We next estimate the blocks of $\mathbf L$.

\paragraph{Cancellation in the overlap kernel.}
\label{ten:frames:part:74}
Fix $U,V$, write $L=L_{U,V}$, and put
\[
 \ell=j-|U\cup V|,\qquad z=j/n.
\]
We claim
\begin{equation}
 \|L\|_\infty\le e^{O(j)} z^{\ell/4},\qquad
 \|L\|_2^2\le e^{O(j)} h_R r_R h_C r_C\, z^{j-\ell}.
 \label{ten:frames:eq:8}
\end{equation}
In the expansion of \(Z\) each term is an expectation using two uniform board tuples \(x,y\). At a prescribed subset of indices they coincide, elsewhere they are independently sampled (different indices always independent across these pairs). The row function in \eqref{ten:frames:eq:7} is evaluated using \(x\), the analogous column function using \(y\), and distinctness factors \(D(x)D(y)\) are included.

For the first inequality, fix the expansion choices and point values on \(U\cup V\). On the remaining coordinates the observed variables are only the row of \(x\) and the column of \(y\), both uniform independently. We couple the two choices in the expansion on each such coordinate: given the observed pair \((r,c)\), we either use \((r,c)\) as both points, or \((r,c'),(r',c)\) respectively with \(c',r'\) independently sampled. All samplings are independent across the remaining coordinates. If some remaining index has no possible distinctness conflict (with other indices, over the choices just described), the signed sum of distinctness factors cancels. The probability that every one of the \(\ell\) indices has a possible conflict is bounded by \(e^{O(j)} z^{\ell/2}\), uniformly in the fixed points. Indeed each vertex has a bounded collection of candidate points, each uniform, collections independent across vertices. Equality edges connect them to fixed points or each other. If they all have conflicts we can choose forest constraints oriented toward roots, some of which may be fixed points, requiring \(h\ge \ell/2\) edges, each specifying which candidates agree. For example, connect vertices in each unanchored connected component along a tree (its size is at least two), and those in anchored components by trees leading to fixed points. The count for given \(h\) is \(\le e^{O(j)}(Cj)^h\) by specifying the parents, with a constant \(C\), and each forest has probability at most $n^{-h}$. To see this, integrate a leaf: the one specified candidate there is uniform and independent of its parent's candidates. Candidates at the same vertex may be correlated, but only one is used when that leaf is removed. Repeating this proves the bound. It is uniform in the fixed marked points and averaged over the observed rows and columns.

Let $p$ be the conditional probability of the conflict event given those observed variables. The absolute averaged signed sum is at most $2^\ell p$, and
\[
 \mathbb E p^2\le\mathbb E p\le e^{O(j)}z^{\ell/2}.
\]
Consequently its kernel from the remaining rows to the remaining columns has Hilbert--Schmidt norm at most $e^{O(j)}z^{\ell/4}$. Finally one integrates the values on \(U\cup V\); their individual marginals on the two sides are uniform so Cauchy--Schwarz suffices. Summation over at most \(2^j\) choices on them proves the operator bound.

For the second inequality we describe diagonal weights for the marked row fibers. Given the full row coordinates and distinct marked positions satisfying the counts \(t_i\), the squared norm of evaluation on the fiber of \(F\) is at most
\begin{equation}
 e^{O(j)} h_R r_R\,d_R(x_U),\qquad
 d_R(x_U)=\operatorname{Tr}\big((P_R)_A Q_A(x_U)\big)/r_R .
 \label{ten:frames:eq:9}
\end{equation}
Here \(Q\) averages over the subgroup of \(H\) fixing these marked positions. This uses transitive evaluation on the marked injections as recalled above, with a conversion of probability normalizations by at most \(e^{O(j)}\). We put \(d_R=0\) if the marked positions do not qualify by themselves. The analogous weight is \(d_C\). They are at most one.

We need the following bound for random marked positions chosen uniformly independently on the board, with any \(u\) of the \(t_R\) positions already prescribed:
\begin{equation}
 \mathbb E\,d_R \ \le\ e^{O(j)} (j/n)^{t_R-u}.
 \label{ten:frames:eq:10}
\end{equation}
To check it, let \(u_i\) be the counts already prescribed in the different rows; assume feasibility. Given row assignments with the right counts and given distinctness of the marked positions, average the fixed-space projection over the remaining choices within rows. In row \(i\), on restricting \(\alpha_i\) to \(\mathfrak S_{m-u_i}\), the only shapes \(\delta\) that can have \(\mathfrak S_{m-t_i}\)-invariants must have first row exactly \(m-t_i\) by branching. The averaged projection on each such block has norm equal to the invariant-space fraction \(\le e^{O(t_i)}\binom{m-u_i}{t_i-u_i}^{-1}\), by \eqref{ten:frames:eq:2}; this averaging is central on \(\mathfrak S_{m-u_i}\). Multiply these bounds and the probability of attaining counts:
\[
 \frac{(t_R-u)!}{b^{t_R-u}\prod_i(t_i-u_i)!}.
\]
The falling factorial bound gives \eqref{ten:frames:eq:10}. This argument covers zero unspecified or zero prescribed positions as well.

For each term of the \(Z\) expansion, the full observed row assignments of \(x\) and column assignments of \(y\) are independent. Consequently we can bound the squared Hilbert--Schmidt norm by integrating over these assignments: within each pair of fibers the operator kernel is a conditional expectation of rank-one forms (evaluations on both sides, times distinctness factors). Jensen bounds its Hilbert--Schmidt norm squared using the products of squared evaluation norms. Hence the factor needed after \(e^{O(j)} h_R r_R h_C r_C\) is bounded by \(\mathbb E[d_R(x_U)d_C(y_V)]\). Let $S\subseteq U\cap V$ be the indices whose marked positions are shared in this expansion term. Conditional on these positions, the other marked points on the two sides are independent. Apply \eqref{ten:frames:eq:10} to the row marks with $S$ prescribed, and then average the column marks without prescriptions. The resulting bound is
\[
 \mathbb E[d_R(x_U)d_C(y_V)]
 \le e^{O(j)}z^{|U|-|S|+|V|}
 \le e^{O(j)}z^{|U\cup V|}.
\]
Triangle inequality over the expansion proves \eqref{ten:frames:eq:8}.

\paragraph{From the blocks to the selected projections.}
Since $0<z<1$, the two block bounds give
\[
 \|L_{U,V}\|_8^8
 \le\|L_{U,V}\|_\infty^6\|L_{U,V}\|_2^2
 \le e^{O(j)}h_Rr_Rh_Cr_C z^{j+\ell/2}
 \le e^{O(j)}h_Rr_Rh_Cr_C z^j.
\]
There are at most $2^j$ mark sets on either side. The triangle inequality for their block embeddings and the bounded right inverses in \eqref{eq:frame-factorization} therefore give the same bound, with a changed absolute constant, for $\|I_R^*I_C\|_8^8$. Its nonzero singular values are those of $P_RP_C$ on $\mathcal I_\lambda$, so
\[
 f^{\lambda_*}\|(P_RP_C)_\lambda\|_8^8
 \le e^{O(j)}h_Rr_Rh_Cr_C(j/n)^j.
\]
Finally $f^\lambda/f^{\lambda_*}\le\binom nj$ and
$\binom nj(j/n)^j\le e^j$ prove \eqref{ten:frames:eq:6}.
\end{proof}

To get \eqref{ten:frames:eq:4}, split \eqref{ten:frames:eq:3} at \(v=8\) into types and dyadic singular value bands on each side. For a term with values \(\le w_R,w_C\) (upper endpoints of the bands), use the selected projections on the inner sides of the two maps and \eqref{ten:frames:eq:6}, multiplying the norm bound by \(w_R w_C\). Within a type, \(\sum w_R(h_R r_R)^{1/8}=O(\log(n!)+1)\) since each term is bounded by \eqref{ten:frames:eq:3}, and an extreme small-value tail is geometric using \(r_R\le h_R\le n!\); likewise for columns. Summing including type counts proves \eqref{ten:frames:eq:4}.

\subsection{The signed-tensor estimate}
\label{ten:frames:part:113}
The selected-rank estimate costs an exponential in the first-row deficit. We now obtain a different cost, depending on boxes outside a small hook. This is a separate signed-tensor argument; it will complement the first estimate when the deficit is large.

Separate the diagram of \(\lambda\) into its first \(s\) rows (shape \(\alpha\)), then the first \(s\) columns of the remaining diagram (transpose shape \(\beta\), i.e., \(\beta\) lists their lengths), then the remaining shape \(\gamma\) of size \(k\). Let \(N'=n-k\) and write \(H_0\) for the entropy in natural log units of the list \((\alpha,\beta)/N'\) (entries are parts). By interleaving tableaux respecting just the indicated rows and columns and the subtableau on \(\gamma\),
\begin{equation}
 f^\lambda\ \le\ f^\gamma \binom nk \exp(N'H_0).
 \label{ten:frames:eq:11}
\end{equation}

Use a mixed tensor space over placements of \(k\) distinct special labels on the board. On the other sites use \(W=\mathbb C^s\oplus\mathbb C^s\), even and odd respectively. Permuting sites uses the usual signed tensor permutation (a minus for interchanging two odd factors); special labels count as even. Explicitly one takes a direct sum of the remaining tensor spaces indexed by the placements and permutes symbols with the sign of the reordering on the odd symbols. This is a representation of \(\mathfrak S_n\). Block diagonal simultaneous unitaries in \(\mathcal G=U(s)\times U(s)\) commute with it. Use the projection \(P\) to \(\mathcal G\)-type \((\alpha,\beta)\). This space contains \(\lambda\) with multiplicity at least \(f^\gamma\). Indeed by ordinary Schur--Weyl applied on each species of tensor factors we have induction of \(\alpha\), the conjugate \(\beta^\intercal\) (because of the sign), and the regular representation for the \(k\) labels. The triple Littlewood--Richardson coefficient with \(\gamma\) is positive: combine \(\beta^\intercal,\gamma\) by row sums, then with \(\alpha\) by column sums (the highest weight sum rule and its conjugate). This yields \(\lambda\).

Write \(Q_H=X^*X,\ Q_K=YY^*\). To bound the trace for the product on \(P\) we use
\[
 \operatorname{Tr}(Q_H P Q_K P).
\]
Since \(P\) preserves placements, only diagonal placement blocks of \(Q_H,Q_K\) occur: in returning each special label one factor preserves its row and the other its column. For any placement the diagonal block \(R_H\) is given by the restriction of the group-algebra coefficients of \(Q_H\) to the subgroup fixing the special positions pointwise. This restriction is positive there (compression in the regular representation). If \(a_i\) special positions occur in row \(i\), its norm on a tuple of types \(\delta_i\vdash m-a_i\) is bounded by
\begin{equation}
 \frac{\exp(O(n^{0.8}))}{\prod_i (m)_{a_i}\, f^{\delta_i}} .
 \label{ten:frames:eq:12}
\end{equation}
Indeed by \eqref{ten:frames:eq:3} with \(v=2\) the full-group identity coefficient times \(|H|\) is at most the type count. On the subgroup the same coefficient times subgroup order bounds dimension times trace on each block. There is an analogous \(R_K\).

We explain a positive-operator majorization for \(R_H\), on this placement. Up to a factor \(\exp(O(n^{0.8}))/\prod_i(m)_{a_i}\), it is majorized by an average of
\begin{equation}
 \bigotimes_{\text{remaining sites }(i,j')} \tau_i ,
 \label{ten:frames:eq:13}
\end{equation}
where the \(\tau_i\) are block diagonal positive matrices on \(W\) of trace one. To verify it use separate block-diagonal unitary types on each row. For \(M=m-a_i\) remaining sites of one row, a type \((u,v')\) with total degree \(M\) corresponds under \(\mathfrak S_M\) to constituents of induction from \(u\) and \((v')^\intercal\), with \(u,v'\) each of length at most \(s\). This again follows by using ordinary Schur--Weyl on the separate species. Signed reorderings used to group factors together do not affect the argument and respect products of even operators. Every such \(\mathfrak S_M\) constituent \(\delta\) satisfies
\[
 u_i\le\delta_i\le u_i+s,\qquad v'_i\le\delta^\intercal_i\le v'_i+s,
\]
with zero padding. For example \((v')^\intercal\) can be obtained by successive induction of at most \(s\) single columns, so Pieri gives the first bound, and the conjugate gives the other. Its hook-length product off the top left \(s\times s\) box is bounded using the horizontal pieces of lengths at most \(u_i\) with hook extensions down by at most \(s\), and the vertical pieces for \(v'_i\) similarly. Thus
\begin{equation}
 f^\delta\ \ge\ \exp\{-O(s^2\log(n+1))\}\ \exp\{M H((u,v')/M)\},
 \label{ten:frames:eq:14}
\end{equation}
by the hook formula and factorial estimates; here \(H(\cdot)\) in an entropy expression denotes Shannon entropy. For instance the hook product is at most \((n+s+1)^{O(s^2)}\prod_i u_i!\,v'_i!\). Empty rows of sites give trivial estimates.

On the other hand take a trace-one matrix of spectrum \((u,v')/M\) in the two blocks and conjugate by Haar block-diagonal unitaries. The expected tensor power of order \(M\) on this type is scalar and at least
\(\exp\{-M H((u,v')/M)-O(s^2\log(n+1))\}\), by the highest weight and dimension bounds. One can average also over type choices at logarithmic cost \(O(s^2\log(n+1))\). Applying this independently in each row gives \eqref{ten:frames:eq:13} from \eqref{ten:frames:eq:12}, \eqref{ten:frames:eq:14}: the accumulated costs fit in \(O(n^{0.8})\) since \(b,m=O(\sqrt n)\). We likewise use a column average of products of \(\sigma_{j'}\).

\subsection{Density normalization on the remaining cells}
\label{ten:frames:part:151}
For these majorizing averages, by positivity it suffices to bound uniformly
\begin{equation}
 \left\|
 (\bigotimes \tau_i^{1/2})\ P\
 (\bigotimes \sigma_{j'}^{1/2})
 \right\|_2^2
 \label{ten:frames:eq:15}
\end{equation}
on the remaining sites of the placement. Use the uniform average over all $b$ rows,
\[
 D_0=\frac1b\sum_{i=1}^b\tau_i+\varepsilon I,\qquad\varepsilon>0.
\]
This averages the full board, including the row states at deleted positions. Insert \((D_0^{-1/2})^{\otimes N'}\) before \((D_0^{1/2})^{\otimes N'} P\) in the middle. The latter operator can be expanded as an integral of the simultaneous \(\mathcal G\) actions, with total absolute coefficient bounded by
\[
 (n+s+1)^{O(s^2)}
 \left((\operatorname{Tr}D_0)^{N'} e^{-N'H_0}\right)^{1/2}.
\]
Indeed this is Fourier expansion on one isotypic type by unitary matrix orthogonality (the extension to \(D_0^{1/2}\) acts there by the same representation with its multiplicity); its operator norm uses the ordered eigenvalues on the two blocks to powers given by the two partitions, bounded by the displayed entropy factor using their sum \(\operatorname{Tr}D_0\). For a simultaneous block unitary \(U\), the remaining tensor product in \eqref{ten:frames:eq:15} after this insertion and expansion has squared Hilbert--Schmidt norm
\[
 \prod_{\text{remaining }(i,j')}
 \operatorname{Tr}(D_0^{-1/2}\tau_i D_0^{-1/2} U \sigma_{j'} U^*) .
\]
Denote these nonnegative entries by $W_{ij'}$ and put $\bar\sigma=m^{-1}\sum_{j'}\sigma_{j'}$. Their full-board average satisfies
\[
 \frac1n\sum_{i,j'}W_{ij'}
 =\operatorname{Tr}\!\left[
 D_0^{-1/2}(D_0-\varepsilon I)D_0^{-1/2}
 U\bar\sigma U^*\right]\le1.
\]
The sum over surviving cells is therefore at most $n$. Arithmetic--geometric mean on those $N'$ cells gives
\[
 \prod_{\text{remaining }(i,j')}W_{ij'}
 \le(n/N')^{N'}\le e^k.
\]
Sending $\varepsilon$ to zero makes $\operatorname{Tr}D_0\to1$ and bounds \eqref{ten:frames:eq:15} by $\exp\{-N'H_0+O(k+n^{0.8})\}$.

It remains to sum the diagonal placement blocks. If $a_i$ and $a'_{j'}$ are their row and column loads, then
\[
 \prod_i(m)_{a_i}\prod_{j'}(b)_{a'_{j'}}
 \ge(m/e)^k(b/e)^k=e^{-2k}n^k.
\]
There are $(n)_k\le n^k$ ordered placements. Thus their number cancels the factor $n^{-k}$ supplied by the two majorizations, at cost $e^{O(k)}$. Combining \eqref{ten:frames:eq:13}, \eqref{ten:frames:eq:15} and the column estimate, we conclude that
\[
 f^\gamma\|(XY)_\lambda\|_2^2 \le \exp\{-N'H_0+O(k+n^{0.8})\}.
\]
This proves \eqref{ten:frames:eq:5} by \eqref{ten:frames:eq:11}.

\subsection{The sparse estimate in short time windows}
\label{ten:frames:part:177}
The two static estimates will close the recursion away from the first row. To finish the argument, we need the following direct estimate for the actual sweep. We retain its time-window proof because it uses a different cancellation from the weighted forest argument.

We record a direct bound, for all sufficiently large \(d\):
\begin{equation}
 \|(\mu_d)_\lambda\|_\infty\ \le\ n^{-c j}
 \quad\text{if }1\le j=n-\lambda_1\le n^{0.9},
 \label{ten:frames:eq:16}
\end{equation}
where \(c>0\) is absolute. Work on \(j\)-point injections \(x,y\) again, with the cube coordinates now and the iid normalization (functions are extended by zero). Any card traveling from \(x_i\) to \(y_i\) in a sweep must use a unique path of sites. For any subset \(I\) of the cards let \(K_I(x,y)\) denote the density on their \(y\) positions relative to uniform independent sampling, given their \(x\) positions. On injections this is a product of pair factors within \(I\): the factor is 1 for non-encounter, and for paths coming to the same switch from opposite sides it is 2 if they split, 0 if they conflict. Such potential encounters for a pair can only start at the step updating their last differing input coordinate. Indeed the unupdated suffixes are fixed. If paths are invalid there is a first failure leading to a zero factor, and on valid paths the formula counts the required fair coins.

In pairing functions of type \(\lambda\) only the part
\[
 \sum_{I\subseteq[j]} (-1)^{j-|I|} K_I
\]
is needed for the full density \(K_{[j]}\), by the top-degree property. This expression vanishes unless every vertex is covered by potential encounters (the paths are determined for all vertices, regardless of validity). We bound the operator norm of each nonnegative \(K_I\) restricted to that event and to injections.

Divide the times into a bounded number of windows updating coordinates $a+1,\ldots,g$, with $g-a\le d/32$. For large $d$ these windows can cover the whole scan. A contact in this window requires the two paths to have the same suffix of $x$ after coordinate $g$ and the same prefix of $y$ through coordinate $a$. Call each such suffix-prefix pair a \emph{window cell}. Coverage of all vertices by contacts therefore implies that, in some window, at least $c_0j$ vertices lie in nonsingleton window cells, for an absolute $c_0>0$.

At least one of $a,d-g$ is at least $0.4d$. We first treat $a\ge0.4d$. Group the input points by their suffix after $g$, and call a group dense when its count is at least
\[
 T=2^a n^{-1/64}.
\]
Split the window event into two parts. Let $\mathcal E_{\rm dense}$ be the event that dense suffix groups contain at least $c_0j/2$ vertices. Let $\mathcal E_{\rm light}$ be window coverage with fewer than that many vertices in dense groups. We will prove, for the mixed sampling associated with each $K_I$,
\begin{equation}\label{eq:frame-window-probabilities}
 \sup_y\Pr_{\rm reverse}(\mathcal E_{\rm dense}\mid y)
 \le e^{Cj}n^{-c_1j},\qquad
 \sup_x\Pr_{\rm forward}(\mathcal E_{\rm light}\mid x)
 \le e^{Cj}n^{-c_1j}.
\end{equation}
These are respectively a column-sum and a row-sum bound for the two positive restricted kernels.

Given $x$, mixed sampling uses genuine sweep paths for $I$ and independent uniform destination bits for the other indices. Reverse sampling is defined in the same way given $y$. Discarding noninjective outputs only decreases the kernel mass. For counts $Z_\nu$ in disjoint site groups at an intermediate or final time, both unrestricted mixed laws satisfy
\begin{equation}
 \mathbb E\prod_\nu(Z_\nu)_{r_\nu}
 \le\prod_\nu(\mathbb E Z_\nu)^{r_\nu}.
 \label{ten:frames:eq:17}
\end{equation}
For the genuine sweep subset, inclusion of any specified distinct sites has probability at most the product of the one-site probabilities. This holds initially and is preserved by random switches: a test set containing one endpoint uses the average of the two old bounds, while a test set containing both uses that their old marginal product is at most the square of their averaged marginal. Summing the site inequalities gives \eqref{ten:frames:eq:17}. Independent extra paths have the same factorial bounds, and independent addition preserves them by the falling-factorial binomial formula.

In reverse sampling from fixed $y$, every suffix group of $x$ has mean $j/2^{d-g}$. The deficit assumption and the window length give
\[
 \frac{j/2^{d-g}}{T}
 =\frac{j}{n}\,2^{g-a}n^{1/64}
 \le n^{0.9-1+1/32+1/64}=n^{-17/320}.
\]
There are at most $j/T$ dense groups. Their possible choices cost at most $\exp(O(j\log n/T))$; if $j<T$, the event is impossible. Since $T\ge n^{0.4-1/64}$, this cost is absorbed by $e^{O(j)}$. The union of any fixed choice of at most $j/T$ groups has mean at most $j n^{-17/320}$. Applying \eqref{ten:frames:eq:17} to a factorial moment of order $\lceil c_0j/2\rceil$ proves the first bound in \eqref{eq:frame-window-probabilities}.

For forward sampling from fixed $x$, each window cell in a nondense suffix group has mean $\mu_\nu\le T/2^a=n^{-1/64}$ at time $a$. On $\mathcal E_{\rm light}$ at least $c_0j/2$ vertices lie in nonsingleton cells of this kind. A cell with $r\ge2$ vertices supplies $\lfloor r/2\rfloor\ge r/3$ disjoint pairs, so there are at least $c_2j$ disjoint co-located pairs, with an absolute $c_2>0$. Since
\[
 \sum_\nu\mu_\nu^2\le n^{-1/64}\sum_\nu\mu_\nu\le j n^{-1/64},
\]
\eqref{ten:frames:eq:17} bounds the expected number of sets of $u=\lceil c_2j\rceil$ such pairs by
\[
 \frac{(\sum_\nu\mu_\nu^2)^u}{2^u u!}
 \le e^{O(j)}n^{-u/64}.
\]
Reducing $c_2$ if necessary handles integer rounding and proves the second bound in \eqref{eq:frame-window-probabilities}.

If instead $d-g\ge0.4d$, transpose the kernel by exchanging $x,y$ and reverse the coordinate names. The window then has endpoints $a'=d-g$, $g'=d-a$, so the preceding argument applies with $a'\ge0.4d$. This exchanges the row and column bounds and leaves the operator norm unchanged.

Each restricted positive kernel thus has one sum bounded by $e^{O(j)}n^{-c_1j}$ and the other by one. The Schur bound gives its operator norm at most $e^{O(j)}n^{-c_1j/2}$. Summing over the bounded number of windows and the $2^j$ subsets $I$ proves \eqref{ten:frames:eq:16}, with a smaller absolute $c$. For $j=1$ the cancelled kernel is zero because there can be no contact.

\subsection{A uniform Schatten exponent}
\label{ten:frames:part:206}
We supply details on uniformity of the Schatten exponent. At any fixed range of sizes a sufficiently large exponent in \eqref{ten:frames:eq:1} suffices. In fact the sweep matrix is a product of orthogonal projections (matching subgroup averages), with norm strictly below one on every nontrivial irreducible: equality would require a joint invariant vector, and all the edges together generate the symmetric group.

Split a large \(d\) into \(\lfloor d/2\rfloor,\lceil d/2\rceil\) bits. Each part of the board update is a product of independent smaller sweeps. If \eqref{ten:frames:eq:1} holds in the smaller problems with exponent \(\le p'\), taking \(p'\ge8\), their independent products on the factor groups have
\[
 h_A\|X_A\|_{p'}^{p'}\le1,\qquad h_B\|Y_B\|_{p'}^{p'}\le1,
\]
by tensoring (norms are contractions). Either product estimate \eqref{ten:frames:eq:4}, \eqref{ten:frames:eq:5} with right hand log \(\mathcal E\) then gives
\begin{equation}
 f^\lambda\|(XY)_\lambda\|_{p'}^{p'} \le \exp(\mathcal E).
 \label{ten:frames:eq:18}
\end{equation}
For clarity this transference uses ordinary Schatten norm complex interpolation: take analytic families in the two group algebras replacing singular values \(x\) by \(x^{z p'/v}\) (zeros held zero), preserving singular vectors. At real part 1 use \eqref{ten:frames:eq:4} or \eqref{ten:frames:eq:5} with its given exponent \(v\), and at real part 0 use product operator norm at most 1. Interpolation at \(v/p'\) gives \eqref{ten:frames:eq:18}, by the three-lines Schatten inequality between operator norm and \(v\)-norm. Thus the log costs do not grow with the starting exponent.

It remains to make one of the two logarithmic costs small relative to $F=\log f^\lambda$. Types of length greater than $n/2$ are annihilated by one matching layer, as noted in Section~\ref{sec:conventions}. For the remaining nontrivial types, Lemma~\ref{lem:surviving-dimension} gives $F\ge c_6j$. We now assume $j\ge n^{0.9}$, leaving the sparse range to \eqref{ten:frames:eq:16}.

With $k,s$ as in \eqref{ten:frames:eq:5}, the hook formula on the lower diagram of size $j$, followed by \eqref{ten:frames:eq:2}, gives
\begin{equation}\label{eq:frames-core-dimension}
 F\ge\log\binom nj+k\log(s/2)-O(j).
\end{equation}
Indeed all hooks of the lower diagram are at most $j$, while the indicated $k$ interior boxes have hooks at most $2j/s$ by their row and column indices there. Comparing their product with $j!$ gives a lower dimension bound $e^{-O(j)}(s/2)^k$.

If $F\ge j d^{1/4}$, choose the selected-rank product estimate \eqref{ten:frames:eq:4}. Its logarithmic cost is
\[
 O(j+n^{0.8})\le O(d^{-1/4}F),
\]
since $j\ge n^{0.9}$.

If $F<j d^{1/4}$, use \eqref{eq:frames-core-dimension} and $F\ge c_6j$ to obtain
\[
 \log(n/j)=O(d^{1/4}),\qquad k=O(F/d).
\]
When $k\ge j/d^2$, we have $\log(n/k)\le\log(n/j)+2\log d=O(d^{1/4}+\log d)$, and hence $k\log(n/k)+k=O(Fd^{-3/4})$. When $0<k<j/d^2$, the crude bound $\log(n/k)\le\log n=O(d)$ gives $k\log(n/k)+k=O(j/d)=O(F/d)$. At $k=0$ these terms vanish. The signed-tensor estimate \eqref{ten:frames:eq:5} therefore has logarithmic cost at most $O(d^{-1/4}F)$ as well; its $n^{0.8}$ term is absorbed using $F\ge c_6n^{0.9}$.

For large enough \(d\), the cost in \eqref{ten:frames:eq:18} is consequently at most \(\varepsilon_d F\) with \(\varepsilon_d=O(d^{-c})\le1/2\), uniformly in these shapes. This also bounds the operator norm to the \(p'\) power by \(\exp(-(1-\varepsilon_d)F)\), so increasing \(p'\) by a factor \(1+O(d^{-c})\) restores \eqref{ten:frames:eq:1}. For $1\le j\le n^{0.9}$, \eqref{ten:frames:eq:16} and $f^\lambda\le n^j$ give, at any sufficiently large absolute exponent $p$,
\[
 f^\lambda\|(\mu_d)_\lambda\|_p^p
 \le(f^\lambda)^2\|(\mu_d)_\lambda\|_\infty^p
 \le n^{(2-cp)j}\le1.
\] The trivial representation satisfies \eqref{ten:frames:eq:1} exactly.

Thus starting with a large enough finite base exponent covering fixed small \(d\) and the direct-bound requirement, it is enough at each large \(d\) to multiply the maximum of the two required smaller exponents by \(1+O(d^{-c})\). These factors have bounded product along bit-length halvings. This proves \eqref{ten:frames:eq:1} with a uniform exponent.

\subsection{Full-deck amplification}
\label{ten:frames:part:231}
Equation~\eqref{ten:frames:eq:1} gives $\|(\mu_d)_\lambda\|_\infty\le(f^\lambda)^{-1/p}$. The Plancherel and dimension-sum argument of Section~\ref{ten:weighted:part:285} therefore applies to repeated forward sweeps, giving worst-start mixing after an absolute number of sweeps, or $O(d)$ physical shuffles. Proposition~\ref{ten:support} supplies the matching lower order. This completes the frame route to the common mixing conclusion.
